\documentclass[11pt,a4paper]{article}

\usepackage[a4paper,margin=2.2cm]{geometry}
\usepackage{fontspec}
\usepackage{polyglossia}
\usepackage{microtype}
\usepackage{booktabs}
\usepackage{tabularx}
\usepackage{array}
\usepackage{xcolor}
\usepackage{enumitem}
\usepackage{hyperref}
\usepackage{titlesec}
\usepackage{fancyhdr}

\setmainlanguage{vietnamese}
\setmainfont{Times New Roman}
\setsansfont{Arial}
\setmonofont{Consolas}

\definecolor{vdtnavy}{HTML}{1F3A5F}
\definecolor{vdtteal}{HTML}{1F7A7A}
\definecolor{vdtgray}{HTML}{F3F5F7}
\definecolor{vdtline}{HTML}{D6DEE6}

\hypersetup{
  colorlinks=true,
  linkcolor=vdtnavy,
  urlcolor=vdtteal,
  citecolor=vdtnavy
}

\pagestyle{fancy}
\fancyhf{}
\lhead{\sffamily Không gian truy xuất bằng chứng}
\rhead{\sffamily \thepage}
\renewcommand{\headrulewidth}{0.4pt}
\setlength{\headheight}{14pt}
\setlength{\emergencystretch}{3em}

\titleformat{\section}
  {\Large\sffamily\bfseries\color{vdtnavy}}
  {\thesection}{0.7em}{}
\titleformat{\subsection}
  {\large\sffamily\bfseries\color{vdtnavy}}
  {\thesubsection}{0.7em}{}

\setlist[itemize]{leftmargin=1.4em,itemsep=0.25em,topsep=0.3em}
\setlist[enumerate]{leftmargin=1.6em,itemsep=0.25em,topsep=0.3em}
\renewcommand{\arraystretch}{1.2}
\newcolumntype{Y}{>{\raggedright\arraybackslash}X}

\newcommand{\method}[1]{\texttt{#1}}
\newcommand{\metric}[1]{\texttt{#1}}
\newcommand{\claim}[1]{\textbf{\color{vdtnavy}#1}}
\newcommand{\bridgecell}{\texttt{tv\_bridge\_title\_}\\\texttt{entities\_rrf}}

\begin{document}

\begin{titlepage}
  \centering
  \vspace*{1.0cm}
  {\Huge\sffamily\bfseries Không gian truy xuất bằng chứng\par}
  \vspace{0.4cm}
  {\Large\sffamily cho QA đa bước và tìm kiếm siêu dữ liệu ngữ nghĩa\par}
  \vspace{1.2cm}
  {\large Báo cáo kỹ thuật nháp cho mentor\par}
  \vspace{0.4cm}
  {\large Cập nhật: 30/06/2026\par}
  \vfill
  \begin{tabularx}{0.86\textwidth}{>{\bfseries\sffamily}l Y}
    Luận điểm &
    Hệ thống tập trung vào tầng truy xuất: tìm đúng bằng chứng, mở rộng lên toàn bộ corpus,
    đánh giá bằng metric phù hợp cho QA đa bước, và làm kết quả truy xuất có thể giải thích được. \\
    Phạm vi &
    Không trình bày như nhật ký sprint. Báo cáo đi theo bài toán, kiến trúc, phương pháp,
    đánh giá, tìm kiếm siêu dữ liệu ngữ nghĩa, giới hạn và bước tiếp theo. \\
  \end{tabularx}
\end{titlepage}

\tableofcontents
\newpage

\section{Tóm tắt}

Dự án xây dựng một không gian truy xuất bằng chứng cho bài toán hỏi đáp đa bước
và tìm kiếm siêu dữ liệu ngữ nghĩa. Trọng tâm của hệ thống không phải sinh câu trả lời cuối cùng,
mà là trả về các tài liệu/ngữ cảnh có thể dùng làm bằng chứng cho bước suy luận phía sau.

Ba điểm chính cần mentor nắm:

\begin{itemize}
  \item \claim{Mở rộng toàn bộ tập dữ liệu:} hệ thống chạy trên toàn bộ tập HotpotQA với
  5,233,329 tài liệu, dùng Elasticsearch cho BM25/kho tài liệu và TurboVec cho truy xuất dense nén.
  \item \claim{Đánh giá đúng bài toán đa bước:} metric quan trọng nhất là
  \metric{full\_support\_recall@10}, vì một truy vấn HotpotQA có thể cần đủ nhiều tài liệu hỗ trợ.
  \item \claim{Tìm kiếm siêu dữ liệu ngữ nghĩa có thể giải thích:} truy vấn tự nhiên có thể được tách thành
  \texttt{effective\_query} và \texttt{metadata\_filters}, nhưng bộ phân tích chỉ chạy khi được bật thủ công để không làm nhiễu truy vấn gốc.
\end{itemize}

Kết quả mới nhất đáng nhấn mạnh là hướng \method{tv\_bridge\_title\_entities\_rrf}.
Trong pilot 200 truy vấn, phương pháp này nâng \metric{full\_support@10} từ 0.545 của
\method{tv\_hybrid} lên 0.620. Cấu hình tuning \texttt{beam1\_terms6} giữ nguyên
0.620 và giảm độ trễ p95 xuống khoảng 1225 ms.

\section{Định nghĩa bài toán}

Bài toán được định nghĩa là \textbf{truy xuất bằng chứng}, không phải sinh câu trả lời.
Với mỗi câu hỏi hoặc truy vấn tự nhiên, hệ thống trả về top-k tài liệu có khả năng làm
bằng chứng. Với HotpotQA, một câu hỏi thường cần nhiều tài liệu hỗ trợ. Nếu hệ thống chỉ
tìm được một tài liệu hỗ trợ nhưng thiếu tài liệu còn lại, bước suy luận phía sau vẫn có thể thất bại.

Vì vậy, ngoài recall@k, hệ thống dùng \metric{full\_support\_recall@k}. Metric này đo tỷ lệ
truy vấn mà top-k kết quả chứa đủ toàn bộ tài liệu hỗ trợ.

\begin{quote}
Trong truy xuất đa bước, câu hỏi không chỉ là top-k có tài liệu liên quan hay không,
mà là top-k có đủ bộ tài liệu hỗ trợ hay không.
\end{quote}

\section{Kiến trúc: tách luồng xử lý offline và luồng truy vấn online}

Kiến trúc nên được trình bày bằng hai luồng riêng. Cách tách này làm rõ phần nào chạy
trước để xây artifact, phần nào chạy tại thời điểm truy vấn để phục vụ người dùng.

\subsection{Luồng ingest/xử lý offline}

Luồng ingest/xử lý chạy trước thời gian chạy. Mục tiêu là biến corpus thô thành các artifact ổn định
có thể tìm kiếm được.

\begin{center}
\fcolorbox{vdtline}{vdtgray}{
\begin{minipage}{0.92\textwidth}
\ttfamily
Tập dữ liệu gốc\\
-> chuẩn hóa tài liệu\\
-> gán numeric\_id ổn định\\
-> sinh siêu dữ liệu tổng hợp\\
-> chia shard JSONL staging\\
-> xây index BM25 trên Elasticsearch\\
-> sinh shard embedding\\
-> xây index dense TurboVec 4-bit\\
-> kiểm tra số lượng tài liệu và artifact
\end{minipage}
}
\end{center}

Các artifact chính:

\begin{tabularx}{\textwidth}{>{\bfseries}p{0.28\textwidth} Y}
\toprule
Artifact & Vai trò \\
\midrule
Shard JSONL staging & Chuẩn hóa corpus, chia shard, giữ \texttt{numeric\_id} ổn định. \\
Index Elasticsearch & Truy xuất lexical bằng BM25, lọc siêu dữ liệu, hydrate tài liệu. \\
Shard embedding & Vector hóa nội dung bằng BGE-small hoặc model theo từng dataset. \\
TurboVec \texttt{.tvim} & Truy xuất dense nén 4-bit cho toàn bộ corpus. \\
Artifact siêu dữ liệu & Siêu dữ liệu tổng hợp để demo truy xuất có nhận biết metadata. \\
\bottomrule
\end{tabularx}

\subsection{Luồng retrieval online}

Luồng retrieval chạy khi người dùng nhập truy vấn. Luồng này không xử lý lại corpus, mà dùng
các index đã xây sẵn.

\begin{center}
\fcolorbox{vdtline}{vdtgray}{
\begin{minipage}{0.92\textwidth}
\ttfamily
Truy vấn người dùng\\
-> bộ phân tích siêu dữ liệu ngữ nghĩa tùy chọn\\
-> effective\_query + metadata\_filters\\
-> truy xuất BM25 từ Elasticsearch\\
-> truy xuất dense từ TurboVec\\
-> hydrate kết quả dense qua Elasticsearch numeric\_id\\
-> gộp hạng bằng RRF hoặc bridge-aware\\
-> overlay support + highlight + chip đã parse\\
-> kết quả trên dashboard React
\end{minipage}
}
\end{center}

Điểm thiết kế quan trọng là Elasticsearch và TurboVec không thay thế nhau. Elasticsearch giữ
vai trò truy xuất lexical, lọc và kho tài liệu. TurboVec giữ vai trò truy xuất dense
nén ở quy mô toàn bộ corpus. Hai backend được nối bằng \texttt{numeric\_id}; tầng ứng dụng
hydrate kết quả dense và gộp hạng bằng RRF hoặc chiến lược bridge-aware.

\section{Các phương pháp retrieval}

\begin{tabularx}{\textwidth}{>{\ttfamily\raggedright\arraybackslash}p{0.30\textwidth} Y}
\toprule
Phương pháp & Vai trò \\
\midrule
es\_bm25 & Mốc so sánh lexical nhanh, mạnh khi truy vấn có độ trùng entity. \\
tv\_dense & Truy xuất ngữ nghĩa bằng TurboVec, hữu ích khi truy vấn ít trùng từ khóa. \\
tv\_hybrid & BM25 + ứng viên dense, gộp hạng bằng Reciprocal Rank Fusion. \\
tv\_filtered\_hybrid & BM25 prefilter + truy xuất dense bị ràng buộc bởi metadata/filter. \\
\bridgecell & Truy xuất bridge-aware cho tài liệu hỗ trợ thứ hai trong QA đa bước. \\
\bottomrule
\end{tabularx}

Kết quả pilot 200 truy vấn dev trên toàn bộ HotpotQA corpus:

\begin{tabularx}{\textwidth}{>{\ttfamily\raggedright\arraybackslash}p{0.30\textwidth} r r Y}
\toprule
Phương pháp & Full-support@10 & Độ trễ p95 & Diễn giải \\
\midrule
es\_bm25 & 0.365 & 359.63 ms & Mốc so sánh lexical nhanh, nhưng thiếu nhiều bộ support đầy đủ. \\
tv\_dense & 0.515 & 868.00 ms & Đường dense cải thiện rõ trên toàn bộ corpus. \\
tv\_hybrid & 0.545 & 2061--3089 ms & Mốc so sánh hybrid ưu tiên chất lượng, tùy lần chạy/cấu hình. \\
\bridgecell & 0.620 & 1224.99 ms & Kết quả pilot tốt nhất sau tuning \texttt{beam1\_terms6}. \\
\bottomrule
\end{tabularx}

Kết luận kỹ thuật: không có phương pháp thắng tuyệt đối trong mọi hoàn cảnh. BM25 tốt về độ trễ,
dense tốt cho độ bền ngữ nghĩa, hybrid cân bằng lexical/dense, còn truy xuất bridge-aware
đánh trực tiếp vào lỗi thiếu tài liệu hỗ trợ trong truy xuất đa bước.

\section{Các thí nghiệm ablation}

Các thí nghiệm ablation trong dự án không chỉ so sánh phương pháp nào cao điểm hơn. Mục tiêu chính
là xác định phần tăng chất lượng đến từ đâu: index lexical, truy xuất dense, gộp hạng RRF, reranking,
sinh ứng viên bridge-aware, hay độ bền trước paraphrase.

\subsection{Thí nghiệm index: BM25 có title-aware}

BM25 có nhận biết title tăng trọng số các trường gần với title/entity như \texttt{title\_exact},
\texttt{lead\_sentence} và \texttt{title\_repeat\_content}. Kết quả cho thấy các metric lexical
tăng nhẹ, nhưng metric chính của HotpotQA không đổi.

\begin{tabularx}{\textwidth}{>{\ttfamily\raggedright\arraybackslash}p{0.24\textwidth} r r r r}
\toprule
Phương pháp & Recall@10 & MRR@10 & nDCG@10 & Full-support@10 \\
\midrule
es\_bm25 & 0.6025 & 0.7108 & 0.5727 & 0.3650 \\
es\_bm25\_title & 0.6050 & 0.7159 & 0.5786 & 0.3650 \\
\bottomrule
\end{tabularx}

Diễn giải: index có nhận biết title là ablation sạch và rẻ, nhưng không giải quyết lỗi
chính của HotpotQA là thiếu tài liệu hỗ trợ thứ hai. Vì vậy phương pháp này được giữ làm
benchmark ablation, không làm mặc định.

\subsection{Thí nghiệm gộp hạng: RRF so với cross-encoder reranker}

Thí nghiệm reranker dùng \texttt{cross-encoder/ms-marco-MiniLM-L-6-v2} với ngân sách 100 ứng viên.
Kết quả cho thấy reranker cải thiện một số rank-sensitive metrics, nhưng không cải thiện
\metric{full\_support@10}.

\begin{tabularx}{\textwidth}{>{\ttfamily\raggedright\arraybackslash}p{0.26\textwidth} r r r r r}
\toprule
Phương pháp & Full-support@10 & Recall@10 & MRR@10 & nDCG@10 & Độ trễ p95 \\
\midrule
tv\_hybrid & 0.5450 & 0.7500 & 0.8691 & 0.7291 & 2061.4370 ms \\
tv\_hybrid\_rerank & 0.5450 & 0.7550 & 0.9268 & 0.7464 & 1304.2816 ms \\
\bottomrule
\end{tabularx}

Dịch chuyển theo cặp trên 200 truy vấn:

\begin{itemize}
  \item Chỉ RRF thành công: 14.
  \item Chỉ reranker thành công: 14.
  \item Cả hai cùng thành công: 95.
  \item Cả hai cùng thất bại: 77.
  \item Net win của reranker: 0.
\end{itemize}

Diễn giải: chưa có bằng chứng đủ mạnh để nói reranker là hướng chính. Nếu full-support không
tăng, vấn đề không chỉ nằm ở việc sắp xếp lại ứng viên. Cần cải thiện bước sinh ứng viên cho
tài liệu hỗ trợ bị thiếu.

\subsection{Thí nghiệm sinh ứng viên: truy xuất bridge-aware}

Truy xuất bridge-aware được thiết kế để đánh vào lỗi lớn nhất: truy vấn tìm được một tài liệu hỗ trợ
nhưng thiếu tài liệu hỗ trợ còn lại. So với \method{tv\_hybrid} chạy một lượt, phương pháp bridge-aware tạo
truy vấn bước hai từ tín hiệu title/entity/lead-sentence của các ứng viên bước đầu.

\begin{tabularx}{\textwidth}{>{\ttfamily\raggedright\arraybackslash}p{0.30\textwidth} r r r r}
\toprule
Phương pháp & Full-support@10 & Recall@10 & nDCG@10 & Độ trễ p95 \\
\midrule
tv\_hybrid & 0.5450 & 0.7500 & 0.7291 & 1146.5764 ms \\
tv\_two\_hop\_bridge\_rrf & 0.5600 & 0.7450 & 0.6999 & 2773.5883 ms \\
\bridgecell & 0.6200 & 0.7850 & 0.7398 & 2670.3591 ms \\
\bottomrule
\end{tabularx}

Chẩn đoán cho thấy phương pháp bridge title/entities chuyển thêm 15/200 truy vấn thành full-support
thành công và giảm lỗi partial-support từ 82 xuống 66. Đây là ablation đầu tiên tạo mức tăng rõ ràng
trên metric chính.

\subsection{Thí nghiệm siêu tham số: tinh chỉnh truy xuất bridge-aware}

Sau khi truy xuất bridge-aware tạo mức tăng, lưới tinh chỉnh kiểm tra liệu beam rộng có thật sự cần thiết
không. Kết quả tốt nhất cho điểm vận hành hiện tại là \texttt{beam1\_terms6}.

\begin{tabularx}{\textwidth}{>{\ttfamily\raggedright\arraybackslash}p{0.24\textwidth} r r r r r}
\toprule
Cấu hình & Beam & Số terms & Full-support@10 & nDCG@10 & Độ trễ p95 \\
\midrule
beam1\_terms6 & 1 & 6 & 0.6200 & 0.7382 & 1224.9911 ms \\
beam2\_terms4 & 2 & 4 & 0.6100 & 0.7430 & 1758.8852 ms \\
beam2\_terms6 & 2 & 6 & 0.6200 & 0.7423 & 2593.4644 ms \\
beam2\_terms8 & 2 & 8 & 0.6200 & 0.7399 & 1827.7998 ms \\
\bottomrule
\end{tabularx}

Diễn giải: beam rộng hơn không tăng full-support trong pilot này. \texttt{beam1\_terms6}
giữ 0.6200 full-support và giảm mạnh độ trễ p95, nên là cấu hình tiếp theo cần được đánh giá lớn hơn.

\subsection{Thí nghiệm độ bền truy vấn: paraphrase}

Benchmark paraphrase gây áp lực lên truy xuất bằng cách đổi truy vấn ở ba mức: mild, strong và lexical-strong.
Mild/strong chủ yếu kiểm tra đảo trật tự câu, còn lexical-strong mới thật sự thay từ nội dung để gây áp lực lên
truy xuất lexical.

\begin{tabularx}{\textwidth}{>{\ttfamily\raggedright\arraybackslash}p{0.26\textwidth} r r r r}
\toprule
Phương pháp & Gốc & Mild & Strong & Lexical strong \\
\midrule
es\_bm25 & 0.365 & 0.365 & 0.375 & 0.340 \\
tv\_dense & 0.515 & 0.515 & 0.515 & 0.495 \\
tv\_hybrid & 0.535 & 0.515 & 0.515 & 0.480 \\
tv\_filtered\_hybrid & 0.430 & 0.435 & 0.440 & 0.395 \\
\bottomrule
\end{tabularx}

Diễn giải: \method{tv\_dense} ổn định nhất khi độ trùng lexical bị phá
($-0.020$ full-support), còn \method{tv\_hybrid} tốt trên truy vấn gốc nhưng giảm mạnh hơn ở
lexical-strong ($-0.055$). Đây là lý do đường dense vẫn cần thiết, dù BM25/hybrid có thể tốt hơn
ở truy vấn gốc.

\subsection{Thí nghiệm họ retrieval: BM25 so với dense và hybrid}

Đây là ablation nền tảng nhất của HotpotQA toàn corpus: chỉ lexical, chỉ dense,
hybrid RRF, và filtered hybrid. Nó trả lời câu hỏi ``có cần đường dense không, và gộp hạng
có đáng giá không?''.

\begin{tabularx}{\textwidth}{>{\ttfamily\raggedright\arraybackslash}p{0.26\textwidth} r r r r}
\toprule
Phương pháp & Recall@10 & MRR@10 & nDCG@10 & Full-support@10 \\
\midrule
es\_bm25 & 0.6025 & 0.7108 & 0.5727 & 0.3650 \\
tv\_dense & 0.7225 & 0.8472 & 0.7082 & 0.5150 \\
tv\_hybrid & 0.7500 & 0.8681 & 0.7286 & 0.5450 \\
tv\_filtered\_hybrid & 0.6800 & 0.8225 & 0.6735 & 0.4550 \\
\bottomrule
\end{tabularx}

Diễn giải: truy xuất dense bằng TurboVec tạo mức tăng lớn so với BM25 trên toàn bộ corpus
($0.3650 \rightarrow 0.5150$ full-support). Hybrid RRF tiếp tục cải thiện lên 0.5450,
nên dense và gộp hạng đều có bằng chứng. Filtered hybrid thấp hơn vì BM25 prefilter có thể loại
mất tài liệu hỗ trợ thứ hai.

\subsection{Thí nghiệm filtered hybrid: hybrid rộng so với hybrid có ràng buộc metadata}

\method{tv\_filtered\_hybrid} là thí nghiệm ablation giữa hybrid rộng ưu tiên chất lượng và đường có ràng buộc cứng
từ allowlist BM25/filter. Đây là phương pháp quan trọng cho tìm kiếm metadata, nhưng không
nên đọc như phương pháp chất lượng mặc định của HotpotQA.

\begin{tabularx}{\textwidth}{>{\ttfamily\raggedright\arraybackslash}p{0.30\textwidth} r r r r}
\toprule
Phương pháp & Full-support@10 & Recall@10 & nDCG@10 & Độ trễ p95 \\
\midrule
tv\_hybrid & 0.5450 & 0.7500 & 0.7286 & 3089.2229 ms \\
tv\_filtered\_hybrid & 0.4550 & 0.6800 & 0.6735 & 1953.6118 ms \\
\bottomrule
\end{tabularx}

Diễn giải: filtered hybrid giảm độ trễ so với hybrid rộng, nhưng mất 0.090 điểm tuyệt đối
full-support. Đánh đổi này hợp lý cho ràng buộc metadata, không hợp lý để thay mặc định
HotpotQA khi mục tiêu là phủ bằng chứng tối đa.

\subsection{Thí nghiệm cross-dataset: VimQA BM25 so với dense và hybrid}

VimQA kiểm tra giả định ``dense/hybrid luôn thắng BM25'' trên proxy truy xuất tiếng Việt.
Dataset này có 3,623 tài liệu và 9,044 truy vấn/qrels.

\small
\begin{tabularx}{\textwidth}{>{\ttfamily\raggedright\arraybackslash}p{0.18\textwidth} r r r r Y}
\toprule
Phương pháp & Recall@10 & MRR@10 & nDCG@10 & Độ trễ p95 & Quyết định \\
\midrule
es\_bm25 & 0.9627 & 0.8606 & 0.8859 & 84.4191 ms & Mặc định tốt nhất. \\
es\_dense & 0.8716 & 0.7272 & 0.7625 & 115.0396 ms & Dense không thắng BM25. \\
es\_hybrid & 0.9644 & 0.8277 & 0.8609 & 206.3031 ms & Recall nhỉnh hơn nhưng xếp hạng/độ trễ kém hơn. \\
\bottomrule
\end{tabularx}
\normalsize

Diễn giải: đặc tính dataset quyết định phương pháp. Trên VimQA, độ trùng lexical cao và mỗi truy vấn
có một gold context, nên BM25 là phương án tốt nhất khi xét cả xếp hạng. Đây là phản ví dụ quan trọng chống
lại việc nói quá về dense/hybrid.

\subsection{Thí nghiệm metadata filter: thu hẹp không gian ứng viên}

Ablation metadata dùng metadata HotpotQA tổng hợp để đo filter làm hẹp không gian ứng viên thế nào.
Đây không phải benchmark chất lượng BEIR, mà là bằng chứng về hành vi sản phẩm/tìm kiếm.

\small
\begin{tabularx}{\textwidth}{>{\raggedright\arraybackslash}p{0.24\textwidth} Y r r}
\toprule
Kịch bản & Bộ lọc & Số tài liệu khớp & Mức thu hẹp \\
\midrule
Chỉ nội dung & không có & 5,233,329 & 0.0000\% \\
Lọc tác giả & author=Nguyen An & 40,886 & 99.2187\% \\
Lọc ngày tạo & created\_at trong tháng 01/2024 & 222,239 & 95.7534\% \\
Lọc ngày sửa & modified\_at 10--20/01 & 60,589 & 98.8422\% \\
Tác giả + ngày tạo & author + created\_at & 1,793 & 99.9657\% \\
\bottomrule
\end{tabularx}
\normalsize

Diễn giải: metadata có cấu trúc giúp corpus 5.23M tài liệu dễ kiểm tra hơn.
Kịch bản tác giả + ngày tạo thu hẹp không gian ứng viên 99.9657\% trước bước xếp hạng theo nội dung.
Khẳng định an toàn: lọc metadata hoạt động như một cơ chế demo/tìm kiếm có kiểm soát. Không khẳng định:
metadata tổng hợp là tác giả thật của HotpotQA hoặc metadata meeting production.

\subsection{Thí nghiệm chiến lược tài liệu: mức tài liệu so với mức đoạn}

EDA về chiến lược tài liệu kiểm tra liệu chia đoạn rộng có nên là bước nâng chất lượng tiếp theo
ở tầng ingest hay không. Kết quả không ủng hộ việc chọn chia đoạn rộng làm bước đầu tiên.

\begin{tabularx}{\textwidth}{>{\bfseries}p{0.22\textwidth} r r r r}
\toprule
Chỉ số & p50 & p90 & p95 & p99 \\
\midrule
Số token nội dung & 39 & 90 & 111 & 160 \\
Số câu/tài liệu & 2 & 5 & 6 & 8 \\
\bottomrule
\end{tabularx}

\vspace{0.5em}

\begin{tabularx}{\textwidth}{>{\bfseries}p{0.32\textwidth} r r}
\toprule
Chiến lược đoạn & Số đơn vị truy xuất ước tính & Tăng so với mức tài liệu \\
\midrule
80 tokens, stride 40 & 6,217,195 & +18.8\% \\
120 tokens, stride 60 & 5,458,362 & +4.3\% \\
160 tokens, stride 80 & 5,290,896 & +1.1\% \\
\bottomrule
\end{tabularx}

Diễn giải: tài liệu HotpotQA vốn đã ngắn, nên index đoạn toàn diện phần lớn tái tạo lại
index ở mức tài liệu nhưng thêm độ phức tạp gom parent. Hướng tốt hơn trước mắt là title/entity
và truy xuất bridge-aware, không phải chia passage rộng.

\subsection{Thí nghiệm backend dense: Elasticsearch dense so với TurboVec và chỉ BM25}

Sprint 3 cũng có ablation về kiến trúc, không chỉ là bảng metric. Các phương án được cân nhắc:

\begin{tabularx}{\textwidth}{>{\bfseries}p{0.30\textwidth} Y}
\toprule
Phương án & Quyết định \\
\midrule
Elasticsearch dense\_vector toàn corpus & Loại cho quy mô local 5.23M tài liệu vì overhead RAM/HNSW quá rủi ro. \\
Truy xuất toàn quy mô chỉ bằng BM25 & Loại làm mục tiêu cuối vì so sánh dense/hybrid là mục tiêu lõi. \\
Lượng tử hóa TurboVec 2-bit trước & Loại ở bước đầu; BGE-small 384 dim với lượng tử hóa 4-bit là mốc so sánh an toàn hơn. \\
TurboVec 4-bit + ES BM25/kho tài liệu & Chọn: index dense độc lập, nén, và join ngược về ES bằng \texttt{numeric\_id}. \\
\bottomrule
\end{tabularx}

Diễn giải: ablation này giải thích vì sao hệ thống có hai backend truy xuất. TurboVec không thay
Elasticsearch; nó thay đường dense vector toàn corpus, còn Elasticsearch vẫn giữ BM25,
lọc và kho hydrate tài liệu.

\subsection{Thí nghiệm quy mô cũ: nano/100k và truy xuất lặp Elasticsearch}

Các thí nghiệm sớm từng so sánh Elasticsearch dense/hybrid và truy xuất lặp ở quy mô nano/100k.
Chúng hữu ích như lịch sử thử nghiệm, nhưng không nên dùng làm khẳng định chính vì thời gian chạy hiện tại
đã chuyển sang BM25 + TurboVec trên toàn bộ corpus.

\begin{tabularx}{\textwidth}{>{\bfseries}p{0.30\textwidth} Y}
\toprule
Thí nghiệm cũ & Cách trình bày \\
\midrule
Nano/100k BM25 so với ES dense/hybrid & Bằng chứng khả thi ban đầu, đã được thay bằng artifact toàn corpus 5.23M tài liệu. \\
ES iterative hybrid/title/sentence/fast & Loại khỏi đường đang chạy trên toàn corpus vì phụ thuộc các trường ES dense cũ. \\
Paraphrase sớm syn020/syn040/syn060 & Đã được thay bằng benchmark 200 truy vấn mild/strong/lexical-strong có validation. \\
\bottomrule
\end{tabularx}

Diễn giải: chỉ nhắc các thí nghiệm này nếu mentor hỏi lịch sử dự án. Với khẳng định chính, dùng các
ablation toàn corpus và các ablation Sprint 4/5 đã được validate ở trên.

\section{Độ trễ và đánh đổi thời gian chạy}

Độ trễ trong báo cáo này là số đo cục bộ trên laptop Windows khoảng 16 GB RAM, không phải
SLA cho môi trường thật. Vai trò của các số độ trễ là giúp giải thích đánh đổi giữa chất lượng,
chiến lược truy xuất và độ phức tạp thời gian chạy.

\begin{tabularx}{\textwidth}{>{\ttfamily\raggedright\arraybackslash}p{0.30\textwidth} Y r Y}
\toprule
Đường chạy & Vai trò chất lượng & Độ trễ p95 & Đánh đổi \\
\midrule
es\_bm25 & Mốc so sánh lexical nhanh & 359.6319 ms & Nhanh nhất, nhưng full-support thấp. \\
tv\_dense & Truy xuất ngữ nghĩa & 868.0033 ms & Cải thiện full-support, thêm chi phí embedding/tìm kiếm vector. \\
tv\_hybrid & Mốc so sánh chất lượng rộng & 2061--3089 ms & Gộp BM25+dense tốt hơn, nhưng tốn hydrate/gộp hạng và đường dense. \\
tv\_filtered\_hybrid & Đường metadata/bộ lọc & 1953.6118 ms & Nhanh hơn hybrid rộng trong pilot, nhưng mất recall/full-support. \\
\bridgecell, beam1\_terms6 & Pilot HotpotQA tốt nhất & 1224.9911 ms & Giữ full-support 0.6200 sau tuning, nhưng phương pháp phức tạp và chưa đưa làm mặc định. \\
VimQA es\_bm25 & Mặc định tiếng Việt & 84.4191 ms & Đặc tính dataset có độ trùng lexical cao, nên BM25 vừa nhanh vừa xếp hạng tốt. \\
\bottomrule
\end{tabularx}

Phân rã độ trễ theo kiến trúc:

\begin{itemize}
  \item \textbf{Đường BM25} chủ yếu tốn truy vấn lexical trên Elasticsearch và lấy tài liệu.
  \item \textbf{Đường dense} thêm embedding cho truy vấn, tìm kiếm vector bằng TurboVec, rồi hydrate bằng Elasticsearch qua \texttt{numeric\_id}.
  \item \textbf{Đường hybrid} chạy cả BM25 và ứng viên dense, sau đó gộp hạng bằng RRF.
  \item \textbf{Đường filtered hybrid} dùng allowlist BM25/bộ lọc để giảm không gian tìm kiếm, nhưng nếu allowlist thiếu tài liệu hỗ trợ thì full-support giảm.
  \item \textbf{Truy xuất bridge-aware} thêm truy xuất bước hai để tìm tài liệu hỗ trợ bị thiếu; tuning giảm beam/terms để giữ gain nhưng hạ p95.
\end{itemize}

Kết luận thời gian chạy: BM25 là đường dự phòng/demo khi cần phản hồi nhanh; \method{tv\_hybrid} là
mốc so sánh chất lượng đơn giản; \method{tv\_bridge\_title\_entities\_rrf} là ứng viên tốt nhất
cho chất lượng HotpotQA sau tuning, nhưng cần benchmark lớn hơn trước khi làm mặc định. Khi trình bày,
không nên khẳng định độ trễ môi trường thật cho bất kỳ phương pháp nào.

\section{Đánh giá đa bước và chẩn đoán lỗi}

HotpotQA cần đánh giá bằng metric phản ánh đủ bộ bằng chứng. Hệ thống dùng:

\begin{itemize}
  \item \metric{recall@10}: top-10 có tài liệu liên quan hay không.
  \item \metric{mrr@10}: tài liệu liên quan đầu tiên đứng cao tới đâu.
  \item \metric{ndcg@10}: chất lượng xếp hạng theo độ liên quan.
  \item \metric{full\_support\_recall@10}: top-10 có đủ toàn bộ tài liệu hỗ trợ hay không.
\end{itemize}

Chẩn đoán xếp hạng chia lỗi thành các nhóm:

\begin{tabularx}{\textwidth}{>{\bfseries}p{0.28\textwidth} Y}
\toprule
Nhóm lỗi & Ý nghĩa \\
\midrule
Thiếu ứng viên & Tài liệu hỗ trợ không vào tập ứng viên. Reranker không thể cứu case này. \\
Chỉ có một phần support & Có một phần support nhưng chưa đủ toàn bộ bộ bằng chứng. \\
Ứng viên xếp thấp & Tài liệu hỗ trợ có trong ứng viên nhưng bị xếp dưới top-k. Đây là nhóm reranker có thể giúp. \\
Thành công & Đủ tài liệu hỗ trợ trong top-k. \\
\bottomrule
\end{tabularx}

Ablation reranker top-100 cho thấy \method{tv\_hybrid\_rerank} không tạo thắng lợi ròng so với RRF
trên pilot 200 truy vấn: cả hai đều đạt \metric{full\_support@10}=0.545, với 14 truy vấn chỉ RRF
thành công và 14 truy vấn chỉ reranker thành công. Vì vậy hướng cải thiện hợp lý hơn hiện tại là
sinh ứng viên/truy xuất tài liệu hỗ trợ thứ hai, dẫn tới \method{tv\_bridge\_title\_entities\_rrf}.

\section{Độ bền truy vấn và thí nghiệm cross-dataset}

\subsection{Độ bền với paraphrase}

Benchmark paraphrase kiểm tra khi độ trùng lexical bị thay đổi. Kết quả profile lexical-strong:

\begin{tabularx}{\textwidth}{>{\ttfamily\raggedright\arraybackslash}p{0.28\textwidth} r r r}
\toprule
Phương pháp & Full-support@10 gốc & Full-support@10 lexical-strong & Chênh lệch \\
\midrule
es\_bm25 & 0.365 & 0.340 & -0.025 \\
tv\_dense & 0.515 & 0.495 & -0.020 \\
tv\_hybrid & 0.535 & 0.480 & -0.055 \\
tv\_filtered\_hybrid & 0.430 & 0.395 & -0.035 \\
\bottomrule
\end{tabularx}

Điểm rút ra: truy xuất dense ổn định hơn khi độ trùng lexical bị phá. Hybrid vẫn tốt trên truy vấn
gốc, nhưng còn khoảng cách về độ bền cần tiếp tục cải thiện bằng gộp hạng, truy xuất bridge-aware hoặc benchmark
sâu hơn.

\subsection{VimQA Tiếng Việt}

VimQA được dùng như proxy truy xuất tiếng Việt:

\begin{itemize}
  \item question -> truy vấn
  \item unique context -> tài liệu
  \item question-context mapping -> qrel
  \item answer -> metadata tham khảo, không phải mục tiêu truy xuất
\end{itemize}

Đặc tính dataset:

\begin{tabular}{lr}
\toprule
Thành phần & Số lượng \\
\midrule
Tài liệu & 3,623 \\
Truy vấn & 9,044 \\
Qrels & 9,044 \\
\bottomrule
\end{tabular}

\vspace{0.5em}

\small
\begin{tabularx}{\textwidth}{>{\ttfamily\raggedright\arraybackslash}p{0.18\textwidth} r r r r Y}
\toprule
Phương pháp & Recall@10 & MRR@10 & nDCG@10 & Độ trễ p95 & Kết luận \\
\midrule
es\_bm25 & 0.9627 & 0.8606 & 0.8859 & 84.42 ms & Mặc định tốt nhất. \\
es\_dense & 0.8716 & 0.7272 & 0.7625 & 115.04 ms & Dense không thắng BM25. \\
es\_hybrid & 0.9644 & 0.8277 & 0.8609 & 206.30 ms & Recall nhỉnh hơn, xếp hạng/độ trễ kém BM25. \\
\bottomrule
\end{tabularx}
\normalsize

Kết luận: không thể giả định dense hoặc hybrid luôn thắng BM25. Phương pháp tốt nhất phụ thuộc đặc tính của dataset.

\section{Tìm kiếm metadata ngữ nghĩa}

Tìm kiếm metadata ngữ nghĩa trả lời feedback rằng người dùng không nên phải nhập truy vấn nội dung và form bộ lọc
riêng biệt. Hệ thống cho phép một truy vấn tự nhiên chứa cả ý định nội dung và ràng buộc metadata.

Ví dụ truy vấn tiếng Việt:

\begin{center}
\fcolorbox{vdtline}{vdtgray}{
\begin{minipage}{0.92\textwidth}
\ttfamily
tài liệu về lịch sử Việt Nam của Nguyen An trước 31/01/2024
\end{minipage}
}
\end{center}

Bộ phân tích tách thành:

\begin{center}
\fcolorbox{vdtline}{vdtgray}{
\begin{minipage}{0.92\textwidth}
\ttfamily
original\_query = tài liệu về lịch sử Việt Nam của Nguyen An trước 31/01/2024\\
effective\_query = lịch sử Việt Nam\\
metadata\_filters.author = Nguyen An\\
metadata\_filters.created\_at\_to = 2024-01-31
\end{minipage}
}
\end{center}

Thiết kế quan trọng:

\begin{itemize}
  \item Bộ phân tích chỉ chạy khi \texttt{semantic\_metadata=true}.
  \item Bộ phân tích rule-based được dùng trước để dễ kiểm soát và test.
  \item \texttt{original\_query} được giữ cho hiển thị/history/support lookup.
  \item \texttt{effective\_query} được dùng cho truy xuất/highlighting.
  \item Bộ lọc nhập tay override bộ lọc đã parse để điều khiển rõ ràng từ UI vẫn là nguồn quyết định cuối.
  \item Metadata không được append vào embedding text để tránh làm nhiễu vector nội dung.
\end{itemize}

Giới hạn cần nói rõ: tìm kiếm metadata ngữ nghĩa hiện là đường đã cài đặt và có đánh giá smoke/protocol,
nhưng chưa đủ lớn để khẳng định bộ phân tích semantic đạt chất lượng sản phẩm thực tế.

\section{Thời gian chạy theo dataset và lớp giải thích}

Thời gian chạy được thiết kế theo dataset: một API/UI phục vụ nhiều dataset, mỗi dataset có profile riêng
về index, phương pháp, phương pháp mặc định, model embedding, metric chính và metadata support.

Dạng endpoint:

\begin{center}
\fcolorbox{vdtline}{vdtgray}{
\begin{minipage}{0.92\textwidth}
\ttfamily
GET  /datasets\\
GET  /datasets/\{dataset\_id\}/stats\\
GET  /datasets/\{dataset\_id\}/queries\\
GET  /datasets/\{dataset\_id\}/benchmarks\\
GET  /datasets/\{dataset\_id\}/embedding-health\\
POST /datasets/\{dataset\_id\}/search
\end{minipage}
}
\end{center}

Lớp giải thích gồm:

\begin{tabularx}{\textwidth}{>{\bfseries}p{0.30\textwidth} Y}
\toprule
Lớp & Mục đích \\
\midrule
Highlight từ khóa truy vấn & Cho thấy kết quả khớp truy vấn ở đoạn nào. \\
Overlay support & Cho thấy kết quả nào là gold support hit. \\
Tóm tắt gold support & Cho biết tài liệu hỗ trợ nào tìm được/chưa tìm được. \\
Chip metadata đã parse & Cho thấy parser hiểu metadata trong truy vấn thế nào. \\
Chẩn đoán xếp hạng & Cho biết lỗi nằm ở sinh ứng viên hay xếp hạng. \\
\bottomrule
\end{tabularx}

\section{Luồng demo đề xuất}

Demo nên minh họa kiến trúc và bằng chứng, không đi theo sprint.

\begin{enumerate}
  \item Mở dashboard, chỉ selector dataset.
  \item Chọn HotpotQA.
  \item Chạy một truy vấn bằng \method{tv\_hybrid} hoặc \method{tv\_bridge\_title\_entities\_rrf}.
  \item Chỉ phần tóm tắt Gold Support và badge Support Hit.
  \item Chỉ các từ khóa truy vấn được highlight trong title/snippet.
  \item Bật chế độ Semantic Metadata.
  \item Chạy truy vấn: \texttt{tài liệu về lịch sử Việt Nam của Nguyen An trước 31/01/2024}.
  \item Chỉ các chip đã parse: Content, Author, Created before.
  \item Chuyển sang VimQA.
  \item Mở trang benchmark và giải thích vì sao BM25 là mặc định cho VimQA.
\end{enumerate}

\section{Khẳng định an toàn, điều không khẳng định và bước tiếp theo}

\subsection{Khẳng định an toàn}

\begin{itemize}
  \item Hệ thống hỗ trợ truy xuất toàn bộ HotpotQA corpus trên 5,233,329 tài liệu.
  \item Hệ thống có BM25, TurboVec dense, hybrid RRF, metadata-filtered hybrid và truy xuất bridge-aware.
  \item \method{tv\_bridge\_title\_entities\_rrf} đạt \metric{full\_support@10}=0.620 trong pilot 200 truy vấn.
  \item Hệ thống có dataset-first API/UI cho HotpotQA và VimQA.
  \item Proxy truy xuất VimQA đã được stage và benchmark trên 9,044 truy vấn.
  \item Đường Parser/API/UI cho semantic metadata đã được cài đặt và bật theo lựa chọn.
\end{itemize}

\subsection{Điều không khẳng định}

\begin{itemize}
  \item Không khẳng định benchmark HotpotQA có thể so trực tiếp với paper.
  \item Không khẳng định chất lượng production cho meeting search.
  \item Không khẳng định độ trễ môi trường thật.
  \item Không khẳng định tìm kiếm metadata ngữ nghĩa đã được đánh giá đầy đủ.
  \item Không khẳng định VimQA là benchmark BEIR-style native.
\end{itemize}

\subsection{Bước tiếp theo}

\begin{enumerate}
  \item Chạy benchmark lớn hơn cho truy xuất bridge-aware trên HotpotQA dev/test.
  \item Tạo benchmark semantic metadata lớn hơn từ metadata tổng hợp.
  \item Định nghĩa schema metadata thật cho meeting search: speaker, participants, timestamp, meeting title, agenda, action items.
  \item Gom lại validation command và kiểm tra số lượng artifact.
  \item Quyết định phương pháp thời gian chạy mặc định sau khi có bằng chứng lớn hơn pilot 200 truy vấn.
\end{enumerate}

\section{Câu Hỏi Mentor Có Thể Hỏi}

\subsection*{Vì sao dùng TurboVec?}

Vì toàn bộ HotpotQA corpus có 5.23M tài liệu. Elasticsearch rất phù hợp cho BM25 và kho tài liệu,
nhưng truy xuất dense vector toàn quy mô có thể nặng trên phần cứng local. TurboVec cho phép truy xuất dense
nén, rồi hydrate kết quả qua Elasticsearch bằng \texttt{numeric\_id}.

\subsection*{Vì sao không chỉ thêm reranker?}

Reranker chỉ giúp nếu tài liệu liên quan đã vào tập ứng viên nhưng bị xếp hạng thấp. Ablation reranker
top-100 cho thấy net full-support win bằng 0 trên pilot 200 truy vấn. Vì vậy hướng hiệu quả hơn hiện tại
là cải thiện sinh ứng viên/truy xuất tài liệu hỗ trợ thứ hai.

\subsection*{Metadata có phải metadata thật không?}

Chưa. Metadata của HotpotQA và VimQA hiện là tổng hợp để demo cơ chế truy xuất có nhận biết metadata. Nếu chuyển
sang meeting search thật, cần schema thật như speaker, participants, timestamp, agenda và action items.

\subsection*{Truy xuất bridge-aware chứng minh điều gì?}

Nó cho thấy lỗi lớn của truy xuất đa bước không chỉ nằm ở xếp hạng. Khi phương pháp biết khai thác tín hiệu title/entity
để tìm tài liệu hỗ trợ thứ hai, \metric{full\_support@10} tăng từ 0.545 lên 0.620 trong pilot.

\appendix
\section{Phụ lục cấu hình}

Phụ lục này ghi các cấu hình đủ để mentor đối chiếu cách hệ thống đang chạy và cách các số trong báo cáo được sinh ra.
Đây là cấu hình project/runtime hiện tại, không phải cấu hình production.

\subsection{Hồ sơ dataset trong runtime}

\small
\begin{tabularx}{\textwidth}{>{\bfseries}p{0.18\textwidth} Y Y}
\toprule
Mục & HotpotQA & VimQA \\
\midrule
Dataset id & \texttt{beir/hotpotqa/dev} & \texttt{vimqa/all} \\
Index runtime & \texttt{hotpotqa\_full\_bm25\_current} & \texttt{vimqa\_all\_dense\_bkai\_current} \\
Phương pháp & \texttt{es\_bm25}, \texttt{tv\_dense}, \texttt{tv\_hybrid}, \texttt{tv\_filtered\_hybrid} & \texttt{es\_bm25}, \texttt{es\_dense}, \texttt{es\_hybrid} \\
Mặc định & \texttt{tv\_hybrid} & \texttt{es\_bm25} \\
Backend dense & \texttt{turbovec} & \texttt{elasticsearch\_dense\_vector} \\
Model embedding & \texttt{BAAI/bge-small-en-v1.5} & \texttt{bkai-foundation-models/vietnamese-bi-encoder} \\
Số chiều vector & 384 & 768 \\
Metric chính & \texttt{full\_support\_recall@10} & \texttt{recall@10} \\
Metadata filter & Có & Có trong profile hiện tại \\
\bottomrule
\end{tabularx}
\normalsize

\subsection{Cấu hình ingest và artifact}

\begin{tabularx}{\textwidth}{>{\bfseries}p{0.30\textwidth} Y}
\toprule
Hạng mục & Cấu hình \\
\midrule
HotpotQA staging & \texttt{artifacts/hotpotqa\_full/staging}, \texttt{docs-per-file=50000}, \texttt{numeric\_id} ổn định. \\
HotpotQA document count & \texttt{5,233,329} documents; validate bằng \texttt{--expected-count 5233329}. \\
HotpotQA BM25 index & Index vật lý \texttt{hotpotqa\_full\_bm25\_v1}, alias runtime \texttt{hotpotqa\_full\_bm25\_current}. \\
HotpotQA ingest batch & \texttt{--batch-size 256}, progress markers trong \texttt{artifacts/hotpotqa\_full/progress}. \\
TurboVec artifact & \texttt{artifacts/hotpotqa\_full/turbovec/hotpotqa\_bge\_small\_4bit.tvim}. \\
TurboVec vector config & BGE-small, \texttt{dim=384}, \texttt{bit\_width=4}. \\
VimQA staging & \texttt{artifacts/vimqa/all/staging/docs-00000.jsonl}, \texttt{3,623} tài liệu. \\
VimQA queries/qrels & \texttt{evaluation/results/vimqa/vimqa\_queries.tsv} và \texttt{vimqa\_qrels.tsv}, \texttt{9,044} truy vấn. \\
VimQA index & \texttt{vimqa\_all\_bm25\_current} cho BM25 và \texttt{vimqa\_all\_dense\_bkai\_current} cho dense/hybrid. \\
\bottomrule
\end{tabularx}

\subsection{Cấu hình retrieval và benchmark}

\small
\begin{tabularx}{\textwidth}{>{\ttfamily\raggedright\arraybackslash}p{0.26\textwidth} Y}
\toprule
Phương pháp & Cấu hình chính \\
\midrule
es\_bm25 & Elasticsearch \texttt{multi\_match} trên \texttt{title\^{}2} và \texttt{content}; \texttt{top\_k=10} trong benchmark chính. \\
tv\_dense & Encode query bằng BGE-small qua embedding service, search TurboVec full index, hydrate tài liệu bằng Elasticsearch qua \texttt{numeric\_id}. \\
tv\_hybrid & BM25 candidates + TurboVec dense candidates, gộp bằng RRF; benchmark chính dùng \texttt{candidate\_k=100}, \texttt{num\_candidates=100}, \texttt{rrf\_k=30}. \\
tv\_filtered\_hybrid & BM25/filter tạo allowlist \texttt{numeric\_id}; TurboVec search trong allowlist rồi gộp RRF. Dùng cho metadata/filter path, không phải quality default. \\
tv\_hybrid\_rerank & Benchmark-only; rerank top-100 candidates bằng \texttt{cross-encoder/ms-marco-MiniLM-L-6-v2}; không expose làm runtime default. \\
\bridgecell & Sinh truy vấn bước hai từ title/entity/lead-sentence; cấu hình khuyến nghị hiện tại là \texttt{beam1\_terms6}: \texttt{beam\_size=1}, \texttt{max\_bridge\_terms=6}. \\
\bottomrule
\end{tabularx}
\normalsize

Các benchmark HotpotQA trong báo cáo dùng cấu hình pilot:

\begin{center}
\fcolorbox{vdtline}{vdtgray}{
\begin{minipage}{0.92\textwidth}
\footnotesize\ttfamily
dataset = beir/hotpotqa/dev\\
index = hotpotqa\_full\_bm25\_current\\
top\_k = 10\\
max\_queries = 200\\
candidate\_k = 100\\
num\_candidates = 100\\
rrf\_k = 30\\
first\_hop\_k = 5\\
second\_hop\_k = 10\\
context\_chars = 256
\end{minipage}
}
\end{center}

Riêng cấu hình bridge tuning tốt nhất hiện tại:

\begin{center}
\fcolorbox{vdtline}{vdtgray}{
\begin{minipage}{0.92\textwidth}
\footnotesize\ttfamily
method = tv\_bridge\_title\_entities\_rrf\\
config = beam1\_terms6\\
beam\_size = 1\\
max\_bridge\_terms = 6\\
top\_k = 10\\
candidate\_k = 100\\
num\_candidates = 100\\
rrf\_k = 30
\end{minipage}
}
\end{center}

\subsection{Cấu hình runtime, API và cache}

\small
\begin{tabularx}{\textwidth}{>{\ttfamily\raggedright\arraybackslash}p{0.34\textwidth} Y}
\toprule
Biến/endpoint & Giá trị hoặc vai trò \\
\midrule
DATASET\_ID & Mặc định \texttt{beir/hotpotqa/dev}. \\
ELASTICSEARCH\_URL & Mặc định local \texttt{http://localhost:9200}. \\
ELASTICSEARCH\_INDEX & Local default \texttt{hotpotqa\_docs\_current}; Docker/demo dùng \texttt{hotpotqa\_full\_bm25\_current}. \\
EMBEDDING\_SERVICE\_URL & Docker/demo dùng \texttt{http://host.docker.internal:8010/embed}. \\
EMBEDDING\_TIMEOUT\_SECONDS & \texttt{30}. \\
REDIS\_URL & Docker dùng \texttt{redis://redis:6379/0}. \\
SEARCH\_CACHE\_TTL\_SECONDS & Mặc định \texttt{300}; demo warm cache có thể dùng \texttt{86400}. \\
TURBOVEC\_INDEX\_PATH & Trong container: \texttt{/app/artifacts/hotpotqa\_full/turbovec/}\newline \texttt{hotpotqa\_bge\_small\_4bit.tvim}. \\
DEFAULT\_SEARCH\_METHOD & \texttt{tv\_hybrid}. \\
\bottomrule
\end{tabularx}
\normalsize

Các endpoint dataset-first chính:

\begin{center}
\fcolorbox{vdtline}{vdtgray}{
\begin{minipage}{0.92\textwidth}
\footnotesize\ttfamily
GET  /datasets\\
GET  /datasets/\{dataset\_id\}/stats\\
GET  /datasets/\{dataset\_id\}/queries\\
GET  /datasets/\{dataset\_id\}/benchmarks\\
GET  /datasets/\{dataset\_id\}/embedding-health\\
POST /datasets/\{dataset\_id\}/search
\end{minipage}
}
\end{center}

\subsection{Cấu hình semantic metadata}

\begin{tabularx}{\textwidth}{>{\bfseries}p{0.30\textwidth} Y}
\toprule
Mục & Cấu hình/ý nghĩa \\
\midrule
Cờ bật parser & \texttt{semantic\_metadata=true}; nếu không bật, query đi theo đường retrieval bình thường. \\
Parser & Rule-based deterministic parser để tách \texttt{effective\_query} và \texttt{metadata\_filters}. \\
Trường filter hiện có & \texttt{author}, \texttt{created\_at}, \texttt{modified\_at}. \\
Ưu tiên filter & Filter nhập tay từ UI override filter parser sinh ra. \\
Embedding text & Không append metadata vào embedding text; vector chỉ biểu diễn content để tránh nhiễu semantic search. \\
Đường retrieval khi có metadata & \texttt{tv\_hybrid} có filter được route sang \texttt{tv\_filtered\_hybrid}; filter được xem là hard constraint. \\
\bottomrule
\end{tabularx}

\subsection{Lệnh reproduce tối thiểu}

\begin{center}
\fcolorbox{vdtline}{vdtgray}{
\begin{minipage}{0.92\textwidth}
\footnotesize\ttfamily
python scripts/stage\_hotpotqa.py --dataset beir/hotpotqa --output-dir artifacts/hotpotqa\_full/staging --docs-per-file 50000\\
python scripts/es\_hotpotqa.py create-index --index hotpotqa\_full\_bm25\_v1 --alias hotpotqa\_full\_bm25\_current --reset\\
python scripts/es\_hotpotqa.py ingest --index hotpotqa\_full\_bm25\_v1 --staging-dir artifacts/hotpotqa\_full/staging --progress-dir artifacts/hotpotqa\_full/progress --batch-size 256\\
python scripts/es\_hotpotqa.py validate --index hotpotqa\_full\_bm25\_current --expected-count 5233329
\end{minipage}
}
\end{center}

\begin{center}
\fcolorbox{vdtline}{vdtgray}{
\begin{minipage}{0.92\textwidth}
\footnotesize\ttfamily
python -m src.evaluation.benchmark\_es --dataset beir/hotpotqa/dev --index hotpotqa\_full\_bm25\_current --methods es\_bm25,tv\_dense,tv\_hybrid --top-k 10 --max-queries 200 --candidate-k 100 --num-candidates 100 --rrf-k 30 --output evaluation/results/hotpotqa\_full/tv\_full\_200.json --run-dir evaluation/runs/hotpotqa\_full
\end{minipage}
}
\end{center}

\subsection{Cấu hình render báo cáo}

Báo cáo được render bằng MiKTeX/XeLaTeX để hỗ trợ tiếng Việt có dấu:

\begin{center}
\fcolorbox{vdtline}{vdtgray}{
\begin{minipage}{0.92\textwidth}
\footnotesize\ttfamily
latexmk -pvc -view=none -xelatex -interaction=nonstopmode -halt-on-error mentor-report.tex
\end{minipage}
}
\end{center}

\end{document}
